\section*{Bab \ref{ch:b2:metamorphic-facies} --- Reaksi dan Fasies Metamorfik}

\begin{solution}{exo:b2:metamorphic-facies:1}
(a) \omterm{def:b2:metamorphic-facies:metamorphism}{Metamorfisme} kontak, yaitu \omterm{def:b2:metamorphic-facies:metamorphism}{metamorfisme} termal lokal oleh panas intrusi.
(b) \omterm{def:b2:metamorphic-facies:metamorphism}{Metamorfisme} dislokasi, di zona sesar: pengecilan ukuran butir, sebuah milonit.
(c) \omterm{def:b2:metamorphic-facies:metamorphism}{Metamorfisme} dasar samudra, oleh air laut panas di dekat punggungan, dengan rekristalisasi
yang tidak lengkap. (d) \omterm{def:b2:metamorphic-facies:metamorphism}{Metamorfisme} penimbunan, regional, tanpa deformasi atau
skistositas.
\end{solution}

\begin{solution}{exo:b2:metamorphic-facies:2}
$\Delta_r S = 95.40 - 82.80 = \qty{12.60}{J.K^{-1}}$, $\Delta_r V = 4.986 - 4.415 =
\qty{0.571}{J.bar^{-1}}$; kemiringannya $12.60/0.571 = \qty{22.1}{bar/K} = \qty{2.21}{MPa/K}$.
Dalam rentang \qty{100}{\celsius}, tekanan kesetimbangan naik \qty{2.21}{kbar}.
\end{solution}

\begin{solution}{exo:b2:metamorphic-facies:3}
(a) Sekis hijau. (b) Sekis biru. (c) Amfibolit: hornblende dengan plagioklas yang anortitnya lebih
dari \qty{17}{\percent}. (d) Eklogit.
\end{solution}

\begin{solution}{exo:b2:metamorphic-facies:4}
(a) $\rho = M/V$: kianit $162.046/44.15 = \qty{3.670}{g.cm^{-3}}$, andalusit
$162.046/51.52 = 3.145$, silimanit $162.046/49.86 = 3.250$. (b) Kianit: pada tekanan yang lebih
tinggi, suku $V\,(P - 1)$ dalam energi Gibbs bertambah paling sedikit untuk volume yang paling kecil.
(c) Kristal alami mengandung unsur-unsur lain (besi feri, misalnya, menggantikan sebagian
aluminium dalam andalusit dan silimanit alami) dan pengukuran pada spesimen
alami kurang teliti dibandingkan dengan volume molar dalam tabel.
\end{solution}

\begin{solution}{exo:b2:metamorphic-facies:5}
(a) $\Delta_r G = 3900 - 298.15\times8.59 = \qty{+1339}{J} > 0$: kianit. (b) Pada
\qty{1}{bar}, $\Delta_r H = T\Delta_r S$: $T = 3900/8.59 = \qty{454}{K} =
\qty{181}{\celsius}$. (c) $P = (8.59\times673.15 - 3900)/0.737 = 1882/0.737 =
\qty{2554}{bar} = \qty{2.55}{kbar}$, dibandingkan dengan 2.54 yang terukur. (d) Reaksi itu terlalu lambat
untuk berlangsung pada suhu rendah: andalusit bertahan secara metastabil, seperti kebanyakan mineral
bersuhu tinggi di permukaan.
\end{solution}

\begin{solution}{exo:b2:metamorphic-facies:6}
(a) $\Delta_r H = 7700 + 1000\times(121.39 - 122.29) = \qty{6800}{J}$; $\Delta_r S = 296.44
- 285.02 = \qty{11.42}{J.K^{-1}}$; $\Delta_r G = 6800 - 1000\times11.42 = \qty{-4620}{J}$.
(b) Dengan nilai pada \qty{298.15}{K}: $7700 - 1000\times12.60 = \qty{-4900}{J}$.
(c) Silimanit. Metode sederhana itu meleset \qty{280}{J}, sekitar \qty{6}{\percent}; bila dibagi
dengan $\Delta_r V = 0.571$, itu setara dengan sekitar \qty{490}{bar} pada garis, jauh lebih kecil daripada
pengaruh ketidakpastian entalpi.
\end{solution}

\begin{solution}{exo:b2:metamorphic-facies:7}
(a) $\Delta_r G = -1127.4 + 1128.5 = \qty{1.1}{kJ}$ dan $\Delta_r V = (34.15 -
36.93)/10 = \qty{-0.278}{J.bar^{-1}}$: $P = 1100/0.278 = \qty{3957}{bar} \approx
\qty{4.0}{kbar}$. (b) $\Delta_r S = 88.0 - 91.7 = -3.7$; kemiringannya $-3.7/-0.278 =
\qty{13.3}{bar/K}$. (c) Dengan $\Delta_r H = 0$, $P = T\Delta_r S/\Delta_r V = 13.31\,T$:
pada \qty{573.15}{K}, \qty{7629}{bar} $= \qty{7.63}{kbar}$. (d) Aragonit cangkang
metastabil di permukaan: organisme membuatnya, dan aragonit itu bertahan karena hanya berubah menjadi
kalsit secara perlahan.
\end{solution}

\begin{solution}{exo:b2:metamorphic-facies:8}
(a) $-32.4/-1.698 = \qty{19.1}{bar/K}$. (b) $P = (T\Delta_r S - \Delta_r H)/\Delta_r V =
(-32.4\times773.15 + 5000)/(-1.698) = (-25050 + 5000)/(-1.698) = \qty{11808}{bar}
\approx \qty{11.8}{kbar}$. (c) $\sqrt{2.6^2 + 3.6^2 + 1.0^2} = \qty{4.55}{kJ}$ pada
$\Delta_r H$, sehingga $4550/1.698 = \qty{2.7}{kbar}$ pada tekanan. (d) Pada tekanan
fasies eklogit, komponen albit dalam plagioklas tidak stabil dibandingkan dengan
jadeit dan kuarsa; natrium masuk ke piroksen (omfasit), dan plagioklas
menghilang.
\end{solution}

\begin{solution}{exo:b2:metamorphic-facies:9}
(a) Dari 1 ke \qty{2}{kbar}: $1000/34 = \qty{29.4}{bar/K}$; dari 2 ke 3: $1000/25 =
\qty{40.0}{bar/K}$. (b) Dengan $\Delta_r S$ tetap, $\Delta_r V$ berbanding terbalik dengan
kemiringan: volume itu menyusut dengan faktor $40.0/29.4 = 1.36$. (c) Air yang dilepaskan: volume
molarnya turun ketika tekanan naik, sedangkan volume padatan hampir tidak berubah.
\end{solution}

\begin{solution}{exo:b2:metamorphic-facies:10}
(a) Pada \qty{873.15}{K}: kianit $=$ andalusit $(8.59\times873.15 - 3900)/0.737 =
\qty{4885}{bar}$; kianit $=$ silimanit $(12.60\times873.15 - 7700)/0.571 =
\qty{5782}{bar}$; andalusit $=$ silimanit $(4.01\times873.15 - 3800)/(-0.166) =
\qty{1799}{bar}$. (b) Pada \qty{4885}{bar}, di bawah garis kianit $=$ silimanit (5782),
silimanit lebih stabil daripada kianit, dan di atas garis andalusit $=$ silimanit
(1799), lebih stabil daripada andalusit: titik itu terletak di medan silimanit, dan
kesetimbangan kianit $=$ andalusit di sana metastabil. (c) Pada \qty{4}{kbar}: di antara 1799
dan \qty{5782}{bar}, silimanit. Pada \qty{7}{kbar}: di atas \qty{5782}{bar} kianit mengalahkan
silimanit, dan di atas \qty{4885}{bar} kianit mengalahkan andalusit: kianit.
\end{solution}

\begin{solution}{exo:b2:metamorphic-facies:11}
(a) $z = P/(\rho g) = 3.75\times10^8/(2700\times9.81) = \qty{14.2}{km}$, dengan ketidakpastian
$0.25\times10^8/26487 = \qty{0.94}{km}$. (b) $(504 - 10)/14.16 =
\qty{35}{\celsius.km^{-1}}$. (c) Ketiga mineral itu berdampingan dalam kesetimbangan hanya di
\omterm{def:b2:metamorphic-facies:invariant}{titik invarian}: lintasan batuan harus melewatinya, dan reaksi-reaksinya harus
berhenti di sana. Kebanyakan batuan lewat di sampingnya dan mempertahankan satu atau dua polimorf.
\end{solution}

\begin{solution}{exo:b2:metamorphic-facies:12}
(a) X: $T = 10 + 10\times30 = \qty{310}{\celsius}$ (\qty{583.15}{K}); $P = 3000\times
9.81\times30000 = \qty{8.83e8}{Pa} = \qty{8.83}{kbar}$. Y: $T = 10 + 30\times20 =
\qty{610}{\celsius}$ (\qty{883.15}{K}); $P = 2800\times9.81\times20000 = \qty{5.49}{kbar}$.
(b) X: kalsit $=$ aragonit pada $13.31\times583.15 = \qty{7.76}{kbar} < 8.83$, aragonit;
albit $=$ jadeit + kuarsa pada $(32.4\times583.15 - 5000)/1.698 = \qty{8.18}{kbar} <
8.83$, jadeit + kuarsa; kianit $=$ andalusit pada $(8.59\times583.15 - 3900)/0.737 =
\qty{1.51}{kbar}$, di bawah suhu titik tripel: kianit. Y: kalsit $=$
aragonit pada \qty{11.75}{kbar} $> 5.49$, kalsit; albit $=$ jadeit + kuarsa pada
\qty{13.9}{kbar}, albit; di atas titik tripel, kianit $=$ silimanit pada
\qty{6.00}{kbar} $> 5.49$ dan andalusit $=$ silimanit pada \qty{1.56}{kbar} $< 5.49$:
silimanit. (c) X, pada \qty{10}{\celsius\per\km}, adalah lintasan yang lebih dingin: jadeit dan
aragonit pada suhu hanya \qty{310}{\celsius} menempatkannya dalam fasies sekis biru. Y, tiga kali
lebih hangat, berakhir dengan silimanit pada \qty{610}{\celsius} dan tekanan sedang, dalam
fasies amfibolit.
\end{solution}

\begin{solution}{pb:b2:metamorphic-facies:1}
\noindent\textbf{Bagian I.}

\textbf{1.} $\Delta_r H = -2589.9 + 2593.8 = \qty{3.9}{kJ}$; $\Delta_r S = 91.39 - 82.80 =
\qty{8.59}{J.K^{-1}}$; $\Delta_r V = 5.152 - 4.415 = \qty{0.737}{J.bar^{-1}}$.

\textbf{2.} $\qty{7.7}{kJ}$; $\qty{12.60}{J.K^{-1}}$; $\qty{0.571}{J.bar^{-1}}$.

\textbf{3.} $\qty{3.8}{kJ}$; $\qty{4.01}{J.K^{-1}}$; $\qty{-0.166}{J.bar^{-1}}$; masing-masing adalah
yang kedua dikurangi yang pertama, karena andalusit saling meniadakan.

\textbf{4.} Silimanit memiliki entropi terbesar (95.40): mineral itu terletak di sisi suhu tinggi
dari kedua garisnya. Kianit memiliki volume terkecil (4.415): mineral itu terletak di sisi
tekanan tinggi dari kedua garisnya. Andalusit, yang volumenya terbesar, menempati sisi
tekanan rendah.

\textbf{5.} $8.59/0.737 = \qty{11.66}{bar/K}$; $12.60/0.571 = \qty{22.07}{bar/K}$;
$4.01/(-0.166) = \qty{-24.16}{bar/K}$.

\textbf{6.} Silimanit memiliki entropi lebih besar tetapi volume lebih kecil daripada andalusit: $\Delta_r S > 0$
dan $\Delta_r V < 0$. Kenaikan suhu menguntungkan silimanit, begitu pula kenaikan
tekanan, sehingga garis itu bergerak ke tekanan lebih rendah pada suhu lebih tinggi.

\noindent\textbf{Bagian II.}

\textbf{7.} $3900 - 298.15\times8.59 = \qty{+1339}{J}$; $7700 - 298.15\times12.60 =
\qty{+3943}{J}$. Keduanya positif: kianit stabil di permukaan.

\textbf{8.} $T = \Delta_r H/\Delta_r S$: $3900/8.59 = \qty{454.0}{K} = \qty{181}{\celsius}$;
$7700/12.60 = \qty{611.1}{K} = \qty{338}{\celsius}$; $3800/4.01 = \qty{947.6}{K} =
\qty{674}{\celsius}$.

\textbf{9.} Ketika dipanaskan pada \qty{1}{bar}, kianit digantikan andalusit pada
\qty{181}{\celsius}; pada \qty{338}{\celsius} kianit sudah tidak stabil, sehingga kianit $=$
silimanit di sana metastabil; andalusit digantikan silimanit pada
\qty{674}{\celsius}. Yang pertama dan ketiga adalah kesetimbangan stabil.

\textbf{10.} $(8.59\times673.15 - 3900)/0.737 = \qty{2554}{bar} = \qty{2.55}{kbar}$,
\qty{0.01}{kbar} dari hasil percobaan.

\textbf{11.} $(12.60\times973.15 - 7700)/0.571 = 4562/0.571 = \qty{7989}{bar} =
\qty{7.99}{kbar}$, \qty{0.13}{kbar} di atas hasil percobaan.

\textbf{12.} Kianit $=$ andalusit naik dari \qty{181}{\celsius} pada \qty{1}{bar} dengan
kemiringan \qty{11.7}{bar/K}; andalusit $=$ silimanit turun dari sekitar \qty{11.5}{kbar} pada
\qty{200}{\celsius} sampai \qty{1}{bar} pada \qty{674}{\celsius}; kianit $=$ silimanit, yang
lebih curam di antara dua garis yang naik, memotong keduanya di dekat \qty{514}{\celsius} dan \qty{3.9}{kbar}, dan setiap garis hanya stabil pada
satu sisi titik potong itu (diagram dalam bab ini).

\noindent\textbf{Bagian III.}

\textbf{13.} $a_1 = 8.59/0.737 = \qty{11.655}{bar/K}$, $b_1 = 3900/0.737 =
\qty{5291.7}{bar}$; $a_2 = 12.60/0.571 = \qty{22.067}{bar/K}$, $b_2 = 7700/0.571 =
\qty{13485.1}{bar}$.

\textbf{14.} $T = (b_2 - b_1)/(a_2 - a_1) = 8193.4/10.412 = \qty{786.92}{K} =
\qty{513.8}{\celsius}$.

\textbf{15.} $P = 11.655\times786.92 - 5291.7 = \qty{3880}{bar} = \qty{3.88}{kbar}$.

\textbf{16.} $(4.01\times786.92 - 3800)/(-0.166) = (3155.5 - 3800)/(-0.166) =
\qty{3882}{bar}$: tekanan yang sama dalam batas pembulatan data, seperti yang dituntut proposisi dalam bab ini.

\textbf{17.} \qty{514}{\celsius} berjarak \qty{10}{\celsius} dari \qty{504}{\celsius}, di dalam
ketidakpastian \qty{20}{\celsius}; \qty{3.88}{kbar} berjarak \qty{0.13}{kbar} dari 3.75,
di dalam 0.25. Kecocokannya baik.

\textbf{18.} Tidak: \qty{536}{\celsius} berjarak \qty{32}{\celsius} dari hasil percobaan, di luar
ketidakpastiannya. Pengabaian kapasitas kalor bukan galat utama: galat utamanya adalah entalpi.

\noindent\textbf{Bagian IV.}

\textbf{19.} $\Delta_r H$ kianit $\to$ andalusit naik \qty{1000}{J}; pada
$T$ tetap, $P = (T\Delta_r S - \Delta_r H)/\Delta_r V$ turun sebesar $1000/0.737 = \qty{1357}{bar}$.

\textbf{20.} $\Delta T = -1357/(a_2 - a_1) = -1357/10.41 = \qty{-130}{K}$: titik tripel
bergeser ke sekitar \qty{384}{\celsius}.

\textbf{21.} Garis kianit $=$ silimanit turun sebesar $1000/0.571 = \qty{1751}{bar}$, dan
titik tripel bergeser $+1751/10.41 = \qty{+168}{K}$.

\textbf{22.} Setiap entalpi saja, bila diubah $\pm\qty{2.0}{kJ}$, menggeser titik tripel sebesar
$\pm\qty{260}{K}$ (andalusit) atau $\pm\qty{336}{K}$ (silimanit): kalorimetri saja mengizinkan
titik tripel dalam rentang beberapa ratus derajat.

\textbf{23.} Entalpi dalam tabel tidak murni kalorimetrik: dalam batas
ketidakpastiannya, entalpi itu disesuaikan agar garis-garis terhitung cocok dengan percobaan bolak-balik.
Percobaanlah yang menetapkan diagram; tabel membawanya, secara konsisten, ke reaksi-reaksi lain.

\textbf{24.} Titik tripel yang dihitung dari tabel dengan sifat-sifat konstan terletak pada
\textbf{\qty[text-series-to-math=true]{514}{\celsius}}.
\end{solution}
